This chapter provides an overview of disaggregated datacenter architectures and smartNIC-to-storage interactions.  We will first give an overview of
the main concepts and definitions used in the literature, then we will review the main
architectural challenges and solutions proposed by researchers. Finally, we will list some
unresolved challenges that we aim to solve in this thesis.

\section{Overview of concepts and definitions}

Hardware is evolving, notably through the transition from hard disk drives to \acrfull{ssd} devices and the decreasing cost of RAM.
These trends have enabled researchers and companies to explore new approaches to \acrfull{hpc}. Modern \acrshort{ssd}s use the \acrfull{nvme} protocol, a low-latency, high-throughput interface designed for flash storage. \acrshort{nvme} devices can be accessed over the network using \acrfull{nvmeof}, which extends the \acrlong{nvme} command set to support network fabrics (smartNICs, \dots) \cite{noauthor_nvme_2021, noauthor_what_nodate, guz_nvme-over-fabrics_2017}. Such advancements made it possible to design new architecture designs, such as disaggregated datacenters. 

This thesis uses \textit{disaggregation} to denote an architecture in which hardware resources such as compute, memory,
accelerators, and storage are split into pools (or nodes) interconnected through a high-speed network. This design allows
resources to be allocated on demand, improving utilization and availability. Two levels are commonly distinguished.
\textit{Partial disaggregation} separates storage from compute and memory, which remain integrated together; this level has
already been widely implemented since the early 2020s. \textit{Full disaggregation} groups each resource type into dedicated
clusters that communicate over the network. \cite{lin_disaggregated_2020}. Figures \ref{fig:partial-disaggregation} and \ref{fig:full-disaggregation} illustrate these two levels of disaggregation.


\begin{figure}[htbp]
	\centering
	\begin{minipage}{0.48\linewidth}
		\centering
		\includesvg[width=0.8\linewidth]{src/img/partial_disa.svg}
		\caption{Partial disaggregation architecture}
		\label{fig:partial-disaggregation}
	\end{minipage}
	\begin{minipage}{0.48\linewidth}
		\centering
		\includesvg[width=0.8\linewidth]{src/img/full_disa.svg}
		\caption{Full disaggregation architecture}
		\label{fig:full-disaggregation}
	\end{minipage}
\end{figure}

High-speed network communication is a critical aspect of disaggregated architectures. \acrfull{rdma} \cite{rdma_rfc} \cite{rdma_consortium} is a protocol that allows one computer to access memory on another computer without CPU involvement, providing high throughput and low latency. \acrfull{roce} \cite{roce_nvidia_docs} is a variant of \acrshort{rdma} that operates over Ethernet networks, making it more widely compatible with existing infrastructure.


Writing efficient storage applications involves understanding the role of \acrfull{spdk}. \acrshort{spdk} is a set of tools and libraries for building high-performance storage applications, particularly those that interact with \acrshort{nvme} devices. It provides a user-space driver that bypasses the kernel, allowing for lower latency and higher throughput compared to traditional kernel-based drivers \cite{noauthor_storage_nodate}. This development kit will be used in this thesis to implement and evaluate our proposed solutions, as detailed in the next chapters.


\begin{draftbox}
	Sections below (until end of chapter) are draft notes and need to be rewritten and structured.
	TODOs : 
	\begin{itemize}
		\item Rewrite and structure improvements
		\item Table of comparison of solutions
	\end{itemize}
\end{draftbox}


\section{Architectural challenges}

Architectural challenges can be grouped into two categories:

\begin{itemize}
	\item Limitations of traditional monolithic servers, which motivate the transition toward disaggregated architectures.
	\item New challenges introduced by disaggregation, which break assumptions commonly made in earlier systems research and require new design approaches.
\end{itemize}

\subsection{Monolithic challenges}

The main limitation of \acrfull{ma} is \textbf{resource stranding}. When one resource is exhausted (for example, memory), other resources on the same server may remain underutilized, leading to waste.

Another important issue in \acrfull{ma} is availability under failures. A single hardware failure can impact the whole server and make it unavailable.

Finally, hardware replacement and upgrades in \acrfull{ma} often require costly maintenance operations.

\subsection{Disaggregation challenges}

In disaggregated architecture (\acrfull{da}), the network becomes the central communication path between resource clusters. This creates strict latency and bandwidth requirements: \acrshort{cpu}-storage communication typically tolerates millisecond-scale delays, whereas \acrshort{cpu}-memory communication requires nanosecond-scale delays (around 100 ns) \cite{lin_disaggregated_2020}.

Current network stacks also expose strong flexibility-performance trade-offs. \acrshort{rdma} and \acrfull{roce} offer high performance but limited flexibility due to hardware coupling. In addition, \acrshort{rdma} can suffer from head-of-line blocking, congestion, and vendor lock-in because its requirements are not always satisfied by commodity hardware. By contrast, Linux \acrshort{tcp} is more flexible but delivers lower performance and less efficient \acrshort{cpu} utilization \cite{skiadopoulos_high-throughput_2024}.


\section{State-of-the-art designs and solutions} \label{sec:sota}

\subsection{Zero-copy data path}

The term \textit{zero-copy} refers to data transfer between devices without requiring data replication by the
\acrshort{cpu}, thereby avoiding redundant copies and reducing \acrshort{cpu} usage. \authcite{skiadopoulos_high-throughput_2024} designed a new prototype called \textit{ZeroNIC} which can zero-copy data between NIC and accelerators. Its main advantage is that is it agnostic of the accelerator type and can run on \acrshort{cpu}, \acrshort{gpu}, \acrshort{fpga}, …
NIC hardware split header and payload, transferring the latter directly to the receiver applications buffers without intermediates
copies.

Concerning \acrshort{nvme} storage, \authcite{sun_scalio_2025} focused on \acrshort{nvmeof} target offloading which extends zero-copy mechanism,
transferring data to the \acrfull{dpu}'s \acrfull{hca} via peer-to-peer \acrfull{pci} communication (\acrfull{p2pdma}) and thus creating a new
kind of zero-copy data path.

\subsection{Batched write}

Batched write mechanism retains write operations until an arbitrary threshold is achieved, then all these operations are run
consecutively while a group commit ensures atomicity of operations. The goal is to improve throughput and reduce the accelerator
\acrshort{cpu} load. While it is not technically a contribution brought by disaggregation, its evolution into such systems is still relevant
to be mentioned.

\subsection{Split stacks}

Split stacks are networking stacks divided in several planes, each one optimized for a specific task. IO-TCP from
\authcite{kim_rearchitecting_2023} separate the traditional Linux \acrshort{tcp} stack into a stack running on the host \acrshort{cpu}, performing
control operations, and thus called the control plane. Such operations are complex and benefits from features of the \acrfull{cisc} \acrshort{cpu}
(for example branch prediction, vectorized instructions, …). The host stack extends the mTCP \cite{jeong_mtcp_2014} API with methods to offload
file opening, writing, and closing, such that applications have the most minimal changes to do in order to run itself on IO-TCP.

The data plane handles all operations related to data packet preparation and transfer. These operations are simpler in comparison
of the control plane and then fit better the \acrfull{arm} cores of the smartNIC. They also benefit from being stateless which suits the
usage of the smartNIC. File and disk \acrfull{io} are offloaded to otherwise it would interfere with the control plane on host \acrshort{cpu}.
Data plane operations are virtually performed and enable a mechanism of SEND / ECHO which translate in multiple \acrfull{mtu} packets
transfer.

Many difficulties due to the split introduce themselves, like \acrfull{rtt} measurement, retransmission, and error handling, but are
essentially resolved with the set of commands provided by the implementation.

\section{Unresolved challenges}

There are some challenges still left open upon reading papers in \textbf{section \ref{sec:sota}}. This section aims to
list them, with the goal to define a clear approach for working on the master's thesis as well as a precise problem to solve.

IO-TCP \cite{kim_rearchitecting_2023} lacks of dynamic synchronization. This because authors assumes that files don't change
during the content delivery service. In their solution, both host and smartNIC share the same filesystem, but any change
on the host-side \acrfull{fs} isn't propagated to the smartNIC stack.
% -> Distributed FS ?

Another problem is that file reading on \acrshort{nvme} disks is done with direct \acrshort{io} on a regular ext4 file system. \acrshort{spdk} reaches the best
performance, but support for file system isn't very developed for now.

Next, with a smaller amount of connection (typically < 4), Linux \acrshort{tcp} stack outperforms IO-TCP by over 1.5 times.

Batched writes used in Scalio \cite{sun_scalio_2025} to improve throughput of write operations come at the cost of latency
increase. This can lead to non-negligible issues in environment where low latency is critical.

\section{Objectives}

This paper will focus on fully disaggregated datacenter, especially on interaction between \\acrshort{nvme} storage pools and smartNIC.
One of our objectives is to allow SmartNICs to communicate directly with the storage without involving any compute pool.
% TO-DO:  give a more precise definition